\subsection{Application to a real-life log: BPI Challenge 2013}
\label{sec:bpi2013-application}

To assess the applicability of the mapping on real-life data, we derive a
collaborative event log from the \emph{incidents} sub-log of the BPI
Challenge~2013, which records the incident-management process of the Volvo~IT
VINST system~\citep{vandongen2013bpi,steeman2013bpi}. The log comprises 7{,}554
cases and 65{,}533 events. Besides its status (\texttt{concept:name}) and
sub-status (\texttt{lifecycle:transition}), each event carries a set of
organizational attributes, among which \texttt{organization involved}
identifies the IT organizational line (e.g., \emph{Org line C}, \emph{Org line
A2}) responsible for the ticket at a given point in time. It is worth noticing
that an incident is frequently handed over between lines before it is
resolved---the \emph{ping-pong} behaviour analysed by several challenge
submissions---so that a single case interleaves the activity of several
organizational units.

We reinterpret this process collaboratively by treating each organizational
line as a \emph{participant}, i.e., an autonomous pool that runs its own
internal handling process on the ticket while it holds it, and each incident
as a \emph{collaboration instance}. In this way, the projection of a
collaboration instance onto one participant constitutes a \emph{local case},
and the succession of that participant's status changes constitutes its
internal process. We deliberately select \texttt{organization involved}
rather than the finer-grained \texttt{org:group} (support team) or
\texttt{org:resource} (individual), because the organizational line is the
granularity that corresponds to a BPMN pool, whereas teams and individuals
correspond to lanes and resources within a pool. Under this reinterpretation,
the resulting log contains 25~participants and 9{,}846 local cases; 2{,}034
of the 7{,}554 collaboration instances (27.0\%) involve two or more
participants, the remainder being degenerate single-participant
collaborations.

\paragraph{Message detection and send-side synthesis.}
The source system does not record explicit communication events;
inter-participant coordination is only implicit in the change of the
responsible line between consecutive events. We recover it as follows.
Whenever a ticket is transferred to a different line, the entering event is,
in the overwhelming majority of cases, an event with status \texttt{Queued}
(sub-status \emph{Awaiting Assignment}): 4{,}051 of the 4{,}350 line
hand-overs (93.1\%) are surfaced in this way. We therefore read each
\texttt{Queued} event whose participant differs from that of the immediately
preceding event as a \emph{receive observation} of a hand-over, represented
as a \texttt{ReceiveTask} event, exactly as in our first pass over this log.
Unlike that first pass, however, we no longer leave the send side
unrepresented. For each such \texttt{ReceiveTask} we synthesize a
correlated \texttt{SendTask}, attributed to the previous line and time-stamped
identically to the triggering \texttt{Queued} event, reusing the
organizational attributes already recorded on it
(\texttt{org:resource}, \texttt{org:role}, \texttt{org:group}). This
re-attribution is justified empirically: of the 4{,}051 \texttt{Queued}
events surfacing a line change, 98.2\% are executed by the resource of the
\emph{previous} line's handler, i.e., the sender already queues the ticket
into the destination line under its own identity. Synthesizing the send
side is therefore a deterministic re-attribution of data already present on
the observed event, not the fabrication of unobserved data; it is declared,
not hidden, via a residual \texttt{collab:synthesized="true"} attribute set
on the synthesized event only (never on the \texttt{ReceiveTask}, which
remains a direct observation), following our stated policy of flagging
rather than silently correcting or silently enriching upstream data. Both
sides of a pair also share the same message identifier, serialized as a
residual attribute, so that downstream analyses which need a correlated
message instance (rather than the two independent observations that rule M4
of the mapping mints by design, one \texttt{Message} object per send or
receive event) can reconstruct it via the mapping's opt-in correlation
enrichment. The receiving participant is the line stamped on the event
itself (\texttt{collab:participant}), while the sending participant
(\texttt{collab:fromParticipant}) is derived as the participant of the
preceding event, on both the \texttt{ReceiveTask} and its synthesized
\texttt{SendTask}; we therefore treat \texttt{fromParticipant} as a declared
pre-processing inference rather than a recorded fact, exactly as we treat
the synthesized send event itself, and we account for both as a threat to
internal validity (Sect.~\ref{sec:threats}): the send side of every
hand-over in this log is a \emph{derived} observation, not a directly
recorded one, with the 98.2\% empirical agreement bounding, but not
eliminating, the resulting attribution error.

\paragraph{Serialization.}
The collaborative log conforms to the collaborative XES
extension~\citep{delgado2025collab}, using its \texttt{elemType},
\texttt{participant} and \texttt{fromParticipant} attributes; the
collaboration instance is represented as an XES trace, and all original event
attributes are preserved. Of the 69{,}584 events, 61{,}482 are typed as
\texttt{task}, 4{,}051 as \texttt{ReceiveTask}, and 4{,}051 as
\texttt{SendTask} (the synthesized events described above, identifiable via
\texttt{collab:synthesized="true"}). Under mapping rule M4, which mints one
\texttt{Message} object per send or receive event, this yields 8{,}102
\texttt{Message} objects in the resulting OCEL~2.0 log.

\paragraph{Reported quality issues.}
In keeping with our principle of flagging rather than silently repairing
upstream data-quality problems, the converter reports, without correcting:
(i)~1{,}156 \texttt{Queued} events occurring at the entry of a collaboration
instance, which represent transfers from an unobserved origin and are not
minted as messages; (ii)~299 line hand-overs (6.9\%) not surfaced through a
\texttt{Queued} event and consequently not represented as messages under the
stated rule; (iii)~2{,}319 events attributed to the catch-all value
\emph{Other}; (iv)~44 of the 649 support groups appearing under more than one
line, i.e., a violation of the expected group-to-line functional dependency;
(v)~6{,}950 events lacking an \texttt{org:role} value; and (vi)~5 events with
the residual status \emph{Unmatched}. The attributes \texttt{product} and
\texttt{impact} are constant within every case and are thus preserved as
attributes of the collaboration instance.

\begin{table}[t]
\centering
\caption{Structural metrics of the collaborative log derived from the BPI
Challenge 2013 incidents sub-log.}
\label{tab:bpi2013-collab-metrics}
\begin{tabular}{lr}
\toprule
Metric & Value \\
\midrule
Collaboration instances (CI) & 7{,}554 \\
\quad with $\geq 2$ participants & 2{,}034 (27.0\%) \\
\quad with 1 participant (degenerate) & 5{,}520 \\
Events, total & 69{,}584 \\
\quad \texttt{task} & 61{,}482 \\
\quad \texttt{ReceiveTask} (receive observations) & 4{,}051 \\
\quad \texttt{SendTask} (synthesized send observations) & 4{,}051 \\
Distinct participants & 25 \\
Local cases (CI~$\times$~participant pairs) & 9{,}846 \\
Hand-overs represented as a send/receive pair & 4{,}051 \\
\texttt{Message} objects, total (OCEL, one per send/receive event) & 8{,}102 \\
CI with $\geq 1$ message & 1{,}982 \\
Participants per CI (mean / max) & 1.30~/~5 \\
Events per CI (mean / max) & 9.21~/~135 \\
Messages per CI (mean / max) & 0.54~/~25 \\
Local case length (mean / median / max) & 7.07~/~5~/~87 \\
\bottomrule
\end{tabular}
\end{table}
